\chapter{[MED-KIN] \emph{Spezielle Patientengruppe:} Kinder}
\minitoc
%%%%%%%%%%%%%%%%%%%%%%%%%%%%%%%%%%%%%%%%%%%%%%%%%%%%%%%%%%%
%%% Pädiatrie                                           %%%
%%%%%%%%%%%%%%%%%%%%%%%%%%%%%%%%%%%%%%%%%%%%%%%%%%%%%%%%%%%

\section{Prinzipieller Umgang mit Kindern}

\begin{itemize}
    \item schonend
    \item Aufregung vermeiden
    \item Bezugsperson wichtig
    \item Vertrauensperson bzw. Eltern in die Handlung einbeziehen
    \item Kind während Transport am Besten eng bei Vertrauensperson belassen
\end{itemize}

\paragraph{Altersgruppen bei Kinder}
%  TODO MED: PÄD: Altersgrenzen bei Kindern checken (#19)

\cite{LPN3}

\begin{center}
\begin{tabulary}{12cm}{LL}\toprule
\Es{Neugeborenes} & bis zum 28. Lebenstag
\\\addlinespace
\Es{Säugling} & bis zum Ende des 1. Lebensjahr
\\\addlinespace
\Es{Kleinkind} & 1-5 Jahre
\\\addlinespace
\Es{Schulkind} & 6-13 Jahre
\\\addlinespace
\Es{Jugendlicher} & 14-18 Jahre
\\\bottomrule
\end{tabulary}
\end{center}

\Cave{In der Literatur kommt es gelegentlich zu
abweichenden Definitionen. So ist z. B. in den
ERC-Richtlinien zur Reanimation ein Neugeborenes nur bis
unmittelbar nach der Geburt als solches klassifiziert.}

\section{Anatomische und physiologische  Besonderheiten}

Kinder unterscheiden sich wesentlich von erwachsenen Patienten, nicht nur in
ihrer Fähigkeit ihre Umwelt und gesetzte Handlungen zu verstehen, sondern auch
hinsichtlich ihrer Anatomie und Physiologie. 



\Fazit{Kinder sind keine kleinen Erwachsenen!}

\begin{table}[H]
\Captionof{table}{Anatomische und physiologische Besonderheiten von Kindern.}

\begin{tabulary}{12cm}{LL}\addlinespace
\TabSub{Atemfrequenz}
&
Neugeborenes: 40\,\nicefrac{x}{min}

Säugling: 20\,\nicefrac{x}{min}

Kleinkind: 15\,\nicefrac{x}{min}

Beatmung bis zum Säugling immer in  Neutralposition (Kopf nicht
überstrecken)lt. ERC. 

\\\addlinespace
\TabSub{Zeichen der Atemnot}
&
Einziehungen zwischen Rippen

\index{Nasenflügeln}\Es{Nasenflügeln}
beschleunigte Atmung,
Atemgeräusche
\\\addlinespace
\TabSub{Pulstasten beim Säugling}
&
Oberarmarterie (Oberarm Innenseite), Achsel
\\\addlinespace
\TabSub{Herzfrequenz}
&
Neugeborenes: 140\,--\,160\,/\,min.

Säugling: 120\,/\,min.

Kinder: 80-100\,/\,min.

Bradykardie  ist fast immer  hypoxiebedingt $\rightarrow$
lebensgefährlich!

\\\addlinespace
\TabSub{Wasser-, Elektrolyt- und Wärmehaushalt}
&
Kinder haben im Verhältnis zu ihrem  Körpergewicht eine
größere Körperoberfläche  als Erwachsene!

höherer Wasseranteil

Austrocknung und Unterkühlung schneller  möglich!
\\\bottomrule
\end{tabulary}
\end{table}

\section{Allgemeines zu Erkrankungen im Kindes- und Jugendalter}

Grundsätzlich können Kinder viele der bereits erwähnten Erkrankungen auch bekommen. 

Bei Kindernotfällen können als zusätzlicher Stressfaktor für
den Rettungsdienst stark emotional betroffene Angehörige sein. Daher gilt es besonders ruhig und gezielt zu handeln. Sie sollten sich der Verantwortung bewusst sein, es geht um die Rettung eines jungen Menschen mit langer Lebenserwartung.

Wichtig für den Transport ins Krankenhaus ist, alle erkrankte
Kinder bis zum vollendeten 18. Lebensjahr auf eine
Kinderambulanz zu transportieren.


\section{Speziell die Kindheit und Jugend betreffende Erkrankungen}




\subsection{Fremdkörperaspiration}


\ERROR{GABS}{????-??-??}{Fremdkörperaspiration: Erwachsen -- Kind überarbeiten
und zusammenführen}
{
\A{GABS: unklare Aufteilung -- Fremdkörperaspiration Erwachsener -- Kind

Lt. ERC besteht ein Unterschied in der Behandlung nur bei < 1a (keine
Abdomenkompression, stattdessen Thoraxkompression) 

Abschnitt zu respiratorischen Notfällen nehmen und nur Verweis für <1 a auf
Kinderteil?}
}
% TODO Fremdkörperaspiration: Erwachsen -- Kind überarbeiten und zusammenführen
% (# 39)





\Definition{Verlegung der Atemwege durch Fremdkörper wie
z.\,B. Nahrungsmittel, Spielzeug o.ä.}

\MarginNo{\cite{Erc2005Sec6}} % p 102

\Layout{2}{\TxSymptome}
{
% TODO MED: Fließtext fehlt.
\begin{itemize}
    \item Atemnot
    \item plötzlicher starker Husten
    \item Stridor oder Keuchen
    \item Würgen
    \item evtl. Zyanose
\end{itemize}
}{

}{}{}{}




\MassnahmenNG{Spezielle Maßnahmen}{Atemwegsverlegung}{}{

\begin{itemize}
  \item \TxMassVitBes
  \begin{itemize}
    \item \ce{O2}: 15\LMin
    \item Lagerung: aufrecht, Oberkörper hoch,
    Atemhilfsmuskulatur unterstützen
  \end{itemize}
  \item Manöver (lt. ERC) unterschiedlich zwischen leichter und schwerer Atemwegsverlegung
  \begin{itemize}
    \item \Es{leichte Atemwegsverlegung} Kind weint und kann noch auf Fragen verbal Antworten; deutliches Einatmen ist möglich.
    \begin{itemize}
      \item zum Husten auffordern; überwachen; hospitalisieren
    \end{itemize}
    \item \Es{schwere Atemwegsverlegung} unfähig zu sprechen, atmen oder husten
    \begin{itemize}
      \item 5 Schläge auf den Rücken
      \item 5 Heimlich-Manöver (!beim Säugling jedoch Thoraxkompressionen!)
    \end{itemize}
  \end{itemize}
  \item Bei Atemstillstand Reanimation
\end{itemize}

}

\subsection{Laryngitis, Pseudokrupp}
\label{topic:laryngitis}
\label{topic:pseudokrupp}


\Layout{2}{\TxBeschreibung}
{ Bei der Laryngitis (Pseudokrupp\index{Pseudokrupp}\footnote{"`Pseudokrupp"'
bezieht sich auf die Heiserkeit (griech.: pseudo für „unecht“; schottisch: croup,
„Heiserkeit“\cite{wiki:pseudokrupp}). Der echte Krupp bezeichnet die
Kehlkopfentzündung bei Diphtherie.}) handelt es sich um eine Entzündung des Kehlkopfes. Es können auch die oberen Atemwege durch das Anschwellen der Schleimhäute behindert sein (Pseudokrupp/Kruppsyndrom). Diese
Infektion wird durch Viren verursacht. 
}{
% TODO SLIDE fehlt
}{}{}{}


\Layout{0}{\TxSymptome}
{
Leitsymptom ist ein pfeifendes Geräusch beim Einatmen
\E{(inspiratorischer Stridor)} sowie ein trockener, bellender Husten. Es besteht Heiserkeit und
Dyspnoe. Fieber ist eher diskret und steht nicht im Vordergrund. Die Laryngitis
tritt vor allem im Alter zwischen 3 Monaten und 6 Jahren auf.
}{
\begin{itemize}
    \item Pfeifendes Geräusch beim Einatmen
    \Z{(inspiratorischer Stridor)}.
    \item Trockener, bellender Husten.
    \item Heiserkeit.
    \item Dyspnoe.
    \item Kaum Fieber.
    \item Typisches Alter: 3 Monate -- 6 Jahre.
\end{itemize}
}{}{}{}




\MassnahmenNG{Spezielle Maßnahmen}{Laryngitis}{}{

\begin{itemize}
    \item Eltern und Kind beruhigen.
    \item Feuchte Luft einatmen lassen (evtl. Fenster
    öffnen; feuchte Tücher im Raum aufhängen; Dusche im Badezimmer aufdrehen).
\end{itemize}

}

\subsection{Epiglottitis}
\label{topic:epiglottitis}

\Layout{2}{\TxBeschreibung}
{
Bei der Epiglottitis handelt es sich um eine \E{bakterielle Entzündung} und
daraus folgende \Es{Schwellung des Kehldeckels}. Sie kommt eher selten vor, kann aber
eventuell lebensbedrohlich werden. Die Epiglottitis ist eine schwere Erkrankung, mit
richtiger  Behandlung besteht aber meist eine gute Prognose.
}{
% TODO SLIDE fehlt
}{}{}{}



\Layout{0}{\TxSymptome}
{
Die Epiglottitis geht mit hohem Fieber einher. Es ist ein Atemgeräusch bei der
Einatmung wahrnehmbar (\E{inspiratorischer Stridor}). Der Patient hält den
Mund offen und es ist ein Speichelfluss bemerkbar. Der
Patient ist eher ruhig und spricht kaum, er hat Angst. 
}{
\begin{itemize}
    \item Hoch fieberhafte Erkrankung.
    \item Inspiratorischer Stridor.
    \item Mund offen, Speichelfluss, spricht kaum!
    \item Angst!
\end{itemize}
}{}{}{}

\MassnahmenNG{Spezielle Maßnahmen}{Epiglottitis}{}{

\begin{itemize}
    \item \TxMassVit
    \item Luftbefeuchtung
    \item Kinder sitzen lassen!
    \item Keine Manipulation im Mundraum (Spatel)!
\end{itemize}
}

 \subsubsection{Laryngitis vs. Epiglottitis}
 
Die Unterscheidung beider Krankheiten kann auch für den erfahrenen Sanitäter schwer sein. Darum hier eine Gegenüberstellung der Differenzialdiagnosen Laryngitis (Kruppsymdrom) und Epiglottitis

\begin{table}[H]
\Captionof{table}{Gegenüberstellung Laryngitis vs. Epiglottitis}

\begin{tabulary}{12cm}{LL}\addlinespace
\TabHeading{Laryngitis}
&
\TabHeading{Epiglottitis}
\\\midrule
Wenig bedrohlich
&
Lebensgefährlich
\\\addlinespace
Kind ist weinerlich
&
Ruhiges Kind, apathisch
\\\addlinespace
Kann schlucken
&
Speichelfluss
\\\addlinespace
Kaum Fieber
&
Hohes Fieber
\\\addlinespace
Kind schreit; weint evtl.
&
Sehr leise, verfällt zusehends
\\\bottomrule
\end{tabulary}
\end{table}
 
 
\subsection{SIDS -- Sudden Infant Death Syndrome -- Plötzlicher Kindstod}
\index{plötzlicher Kindstod}
\index{SIDS}
\index{Sudden Infant Death Syndrome}
\label{topic:sids}

\Layout{0}{\TxBeschreibung}
{
Der Plötzliche Kindstod  ist eine plötzliche, einschneidende Katastrophe im
Familienleben. Aus \Es{unbekannter Ursache} -- und aus oft völliger Gesundheit
-- hört das Kind plötzlich auf zu atmen und verstirbt. 

Dieses Schicksal ereilt Kinder im ersten bis vierten Lebensmonat, es ist eine
Häufung in den Monaten Jänner bis März zu beobachten.

}{
\begin{itemize}
    \item Plötzlich, meist im Schlaf.
    \item Vorwiegend 1.-4.Lebensmonat.
    \item Gehäuft Jänner bis März.
    \item Bei Frühgeburten gehäuft.
\end{itemize}
}{}{}{}


\Layout{0}{Ursachen}
{
Es wird viel über die möglichen Ursachen spekuliert. \Z{Als mögliche Ursache
wird oft ein unreifes Atemzentrum und eine zentrale Atemregulationsstörung
angeführt}

Letzten Endes ist die genaue \Es{Ursache unbekannt}, eine
verlässliche Vorhersage, welche Kinder betroffen sein werden,
gibt es nicht. Es wurden lediglich Risikogruppen wie
Frühgeburten und Zwillingsgeschwister von betroffenen
Kindern, identifiziert. Auch scheint die Lagerung in
Bauchlage das Risiko zu erhöhen.

Es muss betont werden, dass eine verlässliche Vorhersage nicht möglich ist und
\Es{die Eltern keine Schuld trifft}!

}{
\begin{itemize}
    \item Ursache unbekannt.
    \item Risikogruppen.
    \item Zentrale Atemregulationsstörung \Z{(unreifes 
    Atemzentrum?)}.
    \item Eltern trifft keine Schuld!
\end{itemize}
}{}{}{}


\MassnahmenNG{Spezielle Maßnahmen}{SIDS}{}{

\begin{itemize}
	\item Reanimation.
	\item Betreuung der Angehörigen. Den Eltern muss das Gefühl
	gegeben werden, dass sie alles für ihr Kind getan haben. Ein solcher
	Unglücksfall tritt schicksalshaft auf und wäre nicht abwendbar gewesen. 
	\item In jedem Fall nimmt die Kriminalpolizei Ermittlungen auf. Darauf müssen
	die Eltern vorbereitet werden. 
	\item Betreuung für die Angehörigen organisieren
	(Kriseninterventionsdienst\opt{REGVIE}{\footnote{In Wien ist die
	\E{Akutbetreuung Wien} der Stadt Wien zuständig. Sie kann über die
	Leitstelle angefordert werden.}})
	\item Zwillingsgeschwister unbedingt hospitalisieren.
	\item Einsatznachbesprechung und eigene Psychohygiene. Derartige Einsätze
	können -- auch für erfahrenes Fachpersonal -- sehr belastend sein. 
	\opt{ORGASBFLO}{\item \ASBFLO{} Meldung an Gesamteinsatzleiter, Erkundigung
	nach Verfügbarkeit von Peers (auch wenn Peer-Angebot nicht angenommen wird!). }
\end{itemize}

}

\subsection{Ertrinkungsunfall}\index{Ertrinkungsunfall}

\Definition{\Es{Ertrinken}: Tod in den ersten 24\,h nach
Unfall: Ertrinken}

ansonsten:
\index{Beinahe-Ertrinken}\T{Beinahe-Ertrinken}

Kleinkinder sind besonders gefährdet, da sie meist nicht
schwimmen können und sich oft schnell und unbemerkt der Aufmerksamkeit der Aufsichtsperson
entziehen. Besonders ungesicherte Gartenteiche, Seen, Swimmingpools,
Badeanstalten oder auch Wannenbäder stellen eine Gefahrenquelle im Alltag dar. 

\MassnahmenNG{Spezielle Maßnahmen}{Ertrinkungsfall}{}{
\begin{itemize}
  \item Wenn erforderlich Reanimation mit fünf
  Initialbeatmungen. Die Reanimation kann auch nach
  länger andauerndem Kreislaufstillstand erfolgreich sein, da bei einer Unterkühlung der Stoffwechsel verlangsamt wird.
  \item \TxBeiVit \TxMassVitBes
  \item Bei Aspiration: Wenn der Verdacht auf eine
  Aspiration besteht, ist der Patient zur Beobachtung zu hospitalisieren.
\end{itemize}
}


\Fazit{Nobody is dead until warm and dead!}


\subsection{Krampfanfälle im Kindesalter}

Grundsätzlich gelten auch bei Kindern die Ausführungen im Abschnitt
\ref{topic:zerebrale-krampfanfaelle}. Eine Besonderheit stellt jedoch der
\T{Fieberkrampf} dar.


\subsubsection{Fieberkrampf}
\label{topic:fieberkrampf}

\Layout{2}{\TxBeschreibung}
{
Kleinkinder unter 5~Jahren können bei Fieber sogenannte Fieberkrämpfe
entwickeln. Sie entsprechen einem vom Gehirn ausgehenden Krampfanfall und sind grundsätzlich
genau so zu behandeln. 

Mitunter ist die Neigung des Kindes zu Fieberkrämpfen schon bekannt und es
wurden bereits spezielle Zäpfchen für den Bedarfsfall verschrieben (und bei
Eintreffen eventuell schon von den Eltern gegeben). 
}{

}{}{}{}



\MassnahmenNG{Spezielle Maßnahmen}{Krampfanfälle im Kindesalter}{}{
    
    Die Behandlung erfolgt grundsätzlich wie beim Erwachsenen, siehe
    \ref{topic:zerebrale-krampfanfaelle}.
    
    Die Körpertemperatur ist immer zu messen!
}



\subsection{Kindesmisshandlung}
\label{topic:kindesmisshandlung}

\Definition{Absichtliche psychische oder physische Gewalt gegen Kinder}


Kindesmisshandlung kommt in \Es{allen sozialen Schichten} vor!

\Layout{2}{Hinweiszeichen}
{
\begin{itemize}
  \item Langes Intervall zwischen  Verletzungszeitpunkt und Verständigung der
  Rettung.
  \item Anamnese und Verletzungsmuster  widersprechen einander!
  \item Mehrere Verletzungen verschiedenen Alters.
  \item Erkennbarer, verletzungsverursachender Gegenstand (Schuhabdruck,
  \ldots).
  \item Zeichen der Vernachlässigung.
  \item Verletzungen bzw. Schmerzen in der Genitalregion (Achtung: Es gibt auch andere
  Ursachen als sexueller Missbrauch!).
  \item "`Untypische"' Verletzungen.
  \item Verhaltensstörungen.
\end{itemize}
}{
\begin{itemize}
  \item Langes Intervall Verletzung -- Hilfeersuchen.
  \item Widerspruch Anamnese -- Verletzungsmuster.
  \item Verletzungen verschiedenen Alters.
  \item Erkennbarer, verletzungsverursachender Gegenstand .
  \item Vernachlässigung.
  \item Verletzungen / Schmerzen in Genitalregion (Achtung: auch andere
  Ursachen!).
  \item "`Untypische"' Verletzungen.
  \item Verhaltensstörungen.
\end{itemize}
}{}{}{}



\MassnahmenNG{Spezielle Maßnahmen}{Verdacht auf Kindesmisshandlung}{}{
    Die Aufgaben des Rettungsdienstes bei Verdacht auf Kindesmisshandlung
besteht in:

\begin{enumerate}
  \item Daran denken!
  \item Umstände und Umfeld abklären
  \item Unstimmigkeiten und Auffälligkeiten auf eine neutrale Weise
  \Es{dokumentieren}
  \item Obiges verlässlich bei der Übergabe an die nachbetreuende Einrichtung
  weiterleiten. Evtl. Name des Gesprächpartners bei der  Übergabe dokumentieren.
  \item \Es{Patienten nicht im Stich lassen}: Gelindestes Mittel wählen,
  evtl. Vorwand für die Hospitalisierung suchen, aber wenn ein begründeter
  Verdacht auf eine Misshandlung besteht und die
  Erziehungsberechtigten unkooperativ sind notfalls auch die Exekutive beiziehen.
\end{enumerate}

\paragraph{Vorsicht!}

Jedenfalls zu \underline{unterlassen} sind:

\begin{itemize}
  \item Voreilige Verdachtsäußerungen. \Z{Nicht alles was stinkt ist ein
  Fisch!}
  \item Andeutungen gegenüber den Erziehungsberechtigten oder Dritten.
\end{itemize}

}


